% !TeX root = ../main.tex

Die durchgeführten Experimente unterliegen verschiedenen Einschränkungen, die bei der
Interpretation der Ergebnisse berücksichtigt werden müssen.

Die Experimente wurden lokal auf einer Datenbank ohne Netzwerkkommunikation durchgeführt.
Dadurch werden mögliche Einflüsse von Netzwerklatenzen und Netzwerkfehlern nicht berücksichtigt.
Zudem hängen die gemessenen Ausführungszeiten von der verwendeten Hardware ab, sodass bei einer anderen
Hardwarekonfiguration abweichende Ergebnisse zu erwarten sind.
Eine weitere Einschränkung stellt die verwendete Concurrency dar.
Die Experimente wurden mit fest definierten Concurrency-Werten durchgeführt und bilden keine realitätsnahe, variable
Systemlast ab.
Dadurch kann nicht beurteilt werden, wie sich ein Retry bei anderen Werten verhält.
Die Ergebnisse zeigen zwar einen Rückgang der Erfolgsraten bei steigender Concurrency, erlauben jedoch keine Aussage
darüber, ab welchem Wert dieses Verhalten in einer anderen Systemumgebung auftritt.
Die Größe der Workload und die Anzahl der Iterationen sind auch begrenzt.
Mit einer Workload von 500 und fünf Iterationen konnten für die untersuchte Testumgebung ausreichend stabile
Ergebnisse erzielt werden.
Für deutlich größere Workloads oder längere Testzeiträume lässt sich jedoch nicht ausschließen, dass sich die
beobachteten Verhältnisse verändern.

Darüber hinaus wurden die untersuchten Fehler künstlich erzeugt.
Dadurch lässt sich gezielt untersuchen, wie sich der Retry-Ansatz bei den jeweiligen Fehlerklassen verhält.
Die dabei verwendeten Fehlerraten entsprechen jedoch nicht notwendigerweise realen
Produktivsystemen.
Insbesondere die in den Experimenten auftretenden hohen Fehlerraten sind für ein produktives System nicht
repräsentativ.
Die Ergebnisse zeigen daher primär die Wirkung von Retries unter kontrollierten Bedingungen und nicht
dessen Verhalten unter einer realen Fehlerverteilung.
Eine weitere Einschränkung besteht darin, dass lediglich drei Fehlerklassen untersucht wurden.
Für Deadlocks, Serialization Failures und Lock Timeouts konnten dadurch konkrete Aussagen über das Verhalten des
Retry-Ansatzes getroffen werden.
Daraus lässt sich jedoch nicht ableiten, dass die untersuchten Strategien auch für andere transiente oder
persistente Fehler gleichermaßen geeignet sind.
Die fehlende Auswertung der SQLSTATE-Codes stellt eine weitere Einschränkung dar.
Die Untersuchung betrachtet die Fehlerklassen gezielt, wodurch eine Überprüfung des Codes nicht weiter notwendig
war neben der Überprüfung, ob nur die tatsächlichen Fehler aufgetreten sind.
In einem produktiven System sollte jedoch geprüft werden, ob ein aufgetretener Fehler tatsächlich für einen Retry
geeignet ist.
Andernfalls besteht die Gefahr, dass Transaktionen wiederholt werden, bei denen ein erneuter Ausführungsversuch
keine Aussicht auf Erfolg hat und dadurch lediglich zusätzliche Systemlast verursacht.

Schließlich sind auch die untersuchten Parameter begrenzt.
Es wurden nur eine begrenzte Anzahl an Retry-Strategien, Delays und Retry-Counts betrachtet.
Es kann daher nicht ausgeschlossen werden, dass andere Parameterkombinationen oder Strategien für die untersuchten
Fehlerklassen ein besseres Verhältnis zwischen Erfolgsrate und Ausführungszeit erreichen.
Dazu wurden in der Analyse die Parameter nur einzeln betrachtet, ohne den Einfluss der Parameter aufeinander zu
beachten.